\section{Introduction}
\label{sec:intro}

\subsection{The cost regime of multi-agent LLM workflows}

A modern agent framework is a directed graph of LLM calls. \blackrim, the
open-source system this paper draws from, dispatches across eight
domain-leading ``citizen'' agents (engineering, frontend, infrastructure,
security, product, data, ML, knowledge) and seven cross-cutting ``worker''
agents (architect, builder, reviewer, tester, researcher, writer, judge). Each
agent's frontmatter declares a default model tier --- haiku, sonnet, or opus
in the Anthropic family --- and that default has, until now, been a static
choice made at agent-definition time.

The economic gap between tiers is roughly $5\times$ between adjacent tiers
under public Anthropic pricing (haiku $\to$ sonnet $\to$ opus), or
$25\times$ end-to-end. Production telemetry from \blackrim's own dogfooding
shows that $> 60\%$ of total spend during a typical engineering session is
attributable to opus dispatches, and roughly $40\%$ of those dispatches
fall into task cells where empirical evidence already suggests sonnet would
suffice.\footnote{See \cref{sec:eval} for the methodology.} Naively
downgrading every dispatch to sonnet would be a disaster; the same telemetry
shows that critical-path tasks --- ADR review, threat modelling, IR
orchestration --- regress measurably when downgraded.

The selection problem is thus inherently \emph{asymmetric}. Overspending on
opus when sonnet sufficed wastes money but produces correct work. Downgrading
to haiku when sonnet was needed produces broken work that costs much more to
detect and repair than the saved API spend. A workable advisor must
embed this asymmetry in its decision rule, not as an afterthought.

\subsection{Why existing tools do not fit cleanly}

Three families of online-learning algorithms could in principle govern this
problem, but each has a friction with our deployment regime:

\begin{enumerate}[leftmargin=*]
  \item \textbf{Multi-armed bandits} (UCB, Thompson Sampling, EXP3) assume
        homogeneous reward distributions. Our tiers are not interchangeable
        arms --- they have a known stochastic dominance ordering on
        \emph{quality} ($\text{opus} \succeq \text{sonnet} \succeq \text{haiku}$
        in expectation, on most tasks) and the inverse on \emph{cost}. A
        plain MAB algorithm can take many trials to learn what the dominance
        prior already encodes.

  \item \textbf{Contextual bandits} (LinUCB, Thompson with linear or neural
        payoffs) handle the heterogeneity, but their regret guarantees require
        either a parametric reward model or a substantial number of samples
        per context cell. Our long tail of task shapes --- $\approx 170$ cells
        across $17$ agents --- has many cells with $< 10$ samples per
        week.

  \item \textbf{Reinforcement learning} (POMDPs, hierarchical RL) handles
        sequential structure but introduces credit-assignment problems we do
        not need: each dispatch is a one-shot decision, and the reward signal
        is local to that dispatch.
\end{enumerate}

What our regime actually wants is a \emph{Bayesian} decision rule with
\emph{asymmetric loss}. The Bayesian framing absorbs the offline landscape
prior cleanly; per-cell Beta-Bernoulli posteriors update incrementally on
each dispatch; the asymmetric loss is encoded as a one-sided constraint on
posterior probability of quality preservation. Constrained Thompson Sampling
\cite{wu2016conservative,amani2019linear} then becomes the natural sampling
rule.

\subsection{Methodological stance: data first, theory second}
\label{sec:stance}

This paper takes an empirical-first stance. We propose CC-TS as a credible
candidate algorithm, derive its theoretical guarantees, and deploy it. But
the paper's claims rest on the data, not on the theory: every quantitative
statement in \cref{sec:eval} is backed by a CSV in
\texttt{data/} produced by a publicly-available script in
\texttt{scripts/}. If the data eventually disconfirms a CC-TS modelling
assumption --- if posteriors are non-stationary, if the Beta-Bernoulli
shape is wrong, if the conformal coverage assumption fails on production
data --- the algorithm bends to the data, not the data to the algorithm.
The open-questions list \texttt{OQ-A1}--\texttt{OQ-A10} in
\cite{blackrim_moa1c} pre-registers the empirical questions whose answers
could promote the system from per-cell tabular CC-TS to contextual
Thompson Sampling, discounted TS, or a continuous-quality Gaussian
posterior. We treat that promotion as a feature of the deployment, not a
threat to the contribution.

\subsection{Contributions}

This paper makes four contributions, summarised here and developed in the
sections cited:

\begin{itemize}[leftmargin=*]
  \item \textbf{Problem formulation} (\cref{sec:problem}). We formalise
        adaptive model-tier selection as a constrained Bayesian decision
        problem with explicit per-task quality-loss tolerance
        $\qtol \in [0,1]$ partitioned into four operational tiers
        (Critical $\qtol = 0$, Strict $\qtol < 0.02$, Moderate
        $\qtol < 0.05$, Lenient $\qtol < 0.10$).
  \item \textbf{Algorithm} (\cref{sec:algorithm}). We propose Conservative
        Constrained Thompson Sampling (CC-TS), with three components:
        (a) hierarchical Beta priors with partial pooling across agents and
        across shapes within an agent; (b) per-cell Beta-Bernoulli posteriors
        updated from telemetry; (c) conformal calibration of quality-drop
        intervals. We give a regret bound under standard sub-Gaussian
        reward assumptions and prove the conservative property: no downgrade
        is committed with credible quality risk above $\qtol$.
  \item \textbf{System and deployment} (\cref{sec:system}). We describe the
        end-to-end deployment in \blackrim: the dispatch path, telemetry
        capture, eval triggering, and the operator-facing
        \texttt{gt dispatch \texttt{-{}-}advise} surface. The implementation
        is open-source and reproducible from the artefacts cited.
  \item \textbf{Empirical evaluation} (\cref{sec:eval}). We report a
        $\tothink{30}$-day production deployment study across $17$ agents
        and $\approx 170$ cells, comparing CC-TS against three baselines
        (opus-default, heuristic dispatcher, $\epsilon$-greedy). We measure
        cost per session, quality regression on a held-out eval suite, and
        cell-level convergence rates.
\end{itemize}

The remainder of the paper is structured as follows. \Cref{sec:related}
surveys the relevant bandit, conformal-prediction, and LLM-routing
literature. \Cref{sec:problem} formalises the problem. \Cref{sec:system}
describes the deployment context. \Cref{sec:algorithm} develops CC-TS.
\Cref{sec:eval} reports the empirical study. \Cref{sec:discussion} discusses
limitations and \cref{sec:future} sketches extensions.
